\section{Computational program: from substrate equations to baryon-like states}
\label{sec:computational-program}

The manuscript becomes scientifically useful only if its mathematical fragments can be turned into calculations.
This section therefore states the immediate numerical program suggested by the cleaned-up theory.
The emphasis is not on a final claim of success, but on identifying the most informative next target and the sequence of calculations required to reach it.

\subsection{Why proton- and neutron-like states are the right next target}

At the current stage of the model, a baryon-scale target is better motivated than an electron-like target.
The reason is structural rather than rhetorical.
The cubic lattice naturally provides a three-channel axis-aligned carrier sector,
\begin{equation}
  \Psi=(\psi_x,\psi_y,\psi_z)^\top,
\end{equation}
so the first localization problem that genuinely tests the non-Abelian side of the proposal is a stable \emph{triplet-bound} state.
That makes proton- and neutron-like states the natural next targets:
one can ask whether the substrate supports two nearby, stable, fm-scale three-channel solutions with comparable total energies and distinct long-range holonomy diagnostics.
By contrast, an electron-only program does not fully exploit the triplet structure that the lattice already offers.

\subsection{Physical scales to be matched}

The baryon sector imposes a much smaller physical length scale than the electron sector.
Using CODATA 2022 values, the proton mass energy is $938.27208943\,\mathrm{MeV}$, the neutron mass energy is $939.56542194\,\mathrm{MeV}$, the neutron-proton mass difference is $1.29333251\,\mathrm{MeV}$, the proton Compton wavelength is $1.32140985360\times 10^{-15}\,\mathrm{m}$, the neutron Compton wavelength is $1.31959090382\times 10^{-15}\,\mathrm{m}$, and the proton rms charge radius is $8.4075\times 10^{-16}\,\mathrm{m}$ \citep{Mohr2025CODATA}.
These numbers do not determine the substrate theory by themselves, but they do determine the physical target window.
The numerical problem is therefore to find localized states whose characteristic support is on the order of a femtometer and whose conserved energy lies in the baryon range.

\subsection{Necessary numerical stages}

A credible simulation program should proceed in four stages.
\begin{enumerate}
\item \textbf{Linear spectrum and branch identification.}
For the prestretched cubic lattice, compute the linearized dispersion relation and the polarization content across the Brillouin zone.
This identifies the relevant narrowband triplet sector and determines where approximate threefold transport is actually present.

\item \textbf{Berry/Wilczek--Zee diagnostics of the triplet sector.}
Compute link variables, plaquette phases, and non-Abelian holonomies on the discretized momentum lattice.
This determines whether the mixed triplet sector really behaves as a transported $U(3)$ bundle in the parameter regime of interest.

\item \textbf{Localized bound-state search.}
Replace pure real-time evolution as the primary tool by a constrained nonlinear search for stationary or periodically breathing localized states.
In practice this means solving either a damped energy-minimization problem or a continuation problem with fixed total energy and prescribed triplet content.

\item \textbf{Real-time stability and diagnostic extraction.}
Once candidate localized states are found, evolve them in real time, test robustness against perturbations, and extract measurable diagnostics:
core radius, total conserved energy, branch composition, far-field connection data, and effective interaction with slow test packets.
\end{enumerate}

\subsection{Practical discretization guidance}

The baryon-scale problem requires a lattice spacing significantly smaller than the target core size.
A practical starting condition is therefore
\begin{equation}
  a \ll R_{\mathrm{target}} \sim 1\,\mathrm{fm},
\end{equation}
with a computational domain large enough that the core and the far field can be separated cleanly.
The simulation must also satisfy the relativistic CFL constraint of the microscopic branch, so shrinking the spatial scale forces a corresponding reduction of the time step.
For this reason, an explicit search for stationary states is computationally more attractive than attempting to discover them by long real-time trial-and-error runs.

\subsection{What would count as a meaningful first success}

A meaningful first success would \emph{not} be a claim that QCD has been reproduced.
It would be a much narrower result:
a stable three-channel localized mode, supported by the actual substrate equations, whose size is fm-scale, whose conserved energy is in the baryon range, whose branch content is genuinely triplet-like, and whose far-field geometric diagnostics distinguish at least two nearby solution families.
Such a result would justify the next layer of interpretation.
Failure to find such states after systematic variation of prestretch, coupling stencil, and nonlinear constitutive law would be equally informative, because it would show that the present substrate hypothesis is missing a necessary localization mechanism.
